\chapter{Conclusion}

This thesis used genome-scale metabolic modelling to characterise two barley rhizospheric SynComs assembled under different strategies. These communities were a core community (\textit{Hv}SC1), assembled from co-occurring strains, and a trait-optimised community (\textit{Hv}SC2), composed of strains previously shown to promote plant growth or protection. The aim was to understand how assembly strategy shapes community structure, metabolic exchange, and stability.
Both communities relied heavily on cross-feeding to meet their metabolic requirements, but differed markedly in the structure of this exchange: \textit{Hv}SC1 showed more symmetric, evenly distributed metabolite exchanges, while \textit{Hv}SC2's exchanges were sparser and more asymmetric in flux magnitudes. Consistent with this structural difference, \textit{Hv}SC1 was substantially more robust to single-strain removal than \textit{Hv}SC2, whose growth fluctuated depending on member removal. Drop-out experiments further identified strain 892 (\textit{Bosea sp.} F3-2) as a keystone strain within \textit{Hv}SC1. These results demonstrate pronounced effects of assembly strategy on community behaviour and are in line with expectations based on experimental data and literature. 
\\


As shown here, even with remaining uncertainties e.g. in model reconstruction, constraint-based modelling can identify distinct metabolic behaviour and interaction patterns in plant-associated bacteria. The thesis also discusses limitations and opportunities to overcome these from both the computational and the biological perspective. All in all, this thesis contributes to an improved mechanistic understanding of the influence of assembly strategy on SynCom behaviour and can support the development of stable, plant-beneficial microbial communities for barley and beyond.
